\setcounter{section}{24}

  \zwischenueberschrift{Übungsaufgaben}

\inputaufgabe
{}
{

Es sei $X$ eine
\definitionsverweis {differenzierbare Mannigfaltigkeit}{}{}
mit
\definitionsverweis {abzählbarer Basis der Topologie}{}{}
und sei

\mathdisp {0 \, \stackrel{  }{\longrightarrow} \,  U   \, \stackrel{  }{ \longrightarrow}   \,  V   \, \stackrel{  }{\longrightarrow} \,  W \,\stackrel{  } {\longrightarrow}  \, 0} {  }

eine
\definitionsverweis {kurze exakte Sequenz}{}{}
von
\definitionsverweis {differenzierbaren Vektorbündeln}{}{}
über $X$. Zeige mithilfe von
[[Mannigfaltigkeit/Abzählbare Basis/Überdeckung/Partition/Fakt|Satz 22.10]],
dass es eine Spaltung der Sequenz gibt, also einen Vektorbündelisomorphismus

\mavergleichskette
{\vergleichskette
{ V
}
{ \cong }{ U \oplus W
}
{  }{
}
{    }{
}
{   }{
}
}
{}{}{,}
der mit den Bündelhomomorphismen verträglich ist.

}
{} {}

\inputaufgabe
{}
{

Es sei

\mavergleichskette
{\vergleichskette
{ W
}
{ \subseteq }{ \R^n
}
{  }{
}
{    }{
}
{   }{
}
}
{}{}{}
\definitionsverweis {offen}{}{,}
\maabb {h} { W } { \R
} {}
eine zweifach
\definitionsverweis {stetig differenzierbare Funktion}{}{}
und

\mavergleichskette
{\vergleichskette
{ Y
}
{ = }{ h^{-1}(c)
}
{  }{
}
{    }{
}
{   }{
}
}
{}{}{}
die
\definitionsverweis {Faser}{}{}
zu

\mavergleichskette
{\vergleichskette
{ c
}
{ \in }{ \R
}
{  }{
}
{    }{
}
{   }{
}
}
{}{}{,}
wobei $h$ in jedem Punkt von $Y$
\definitionsverweis {regulär}{}{}
sei. Das erste und das zweite
\definitionsverweis {Tangentialbündel}{}{}
von $Y$ seien als
\definitionsverweis {abgeschlossene Untermannigfaltigkeit}{}{}
in $W \times \R^n$ bzw. in
\mathl{W \times \R^n \times \R^n \times \R^n}{} im Sinne von
[[Differenzierbare Hyperfläche/Zweites Tangentialbündel/Implizite Beschreibung/Bemerkung|Bemerkung 24.4]]
zusammen mit den Bündelprojektionen
\maabbeledisp {\pi} {TY} {Y
} {(P,u)} {P
} {,}
und
\maabbeledisp {q} {TTY} {TY
} {(P,u,v,w)} {(P,u)
} {,}
realisiert. Zeige die folgenden Aussagen.
\aufzaehlungdrei{Eine
\definitionsverweis {differenzierbare Kurve}{}{}
\maabbeledisp {\epsilon} { I } { TY
} { t } { (P(t), u(t))
} {,}
auf einem offenen Intervall

\mavergleichskette
{\vergleichskette
{ I
}
{ \subseteq }{ \R
}
{  }{
}
{    }{
}
{   }{
}
}
{}{}{}
mit

\mavergleichskette
{\vergleichskette
{ 0
}
{ \in }{ I
}
{  }{
}
{    }{
}
{   }{
}
}
{}{}{}
definiert den
\definitionsverweis {Tangentialvektor}{}{}

\mavergleichskette
{\vergleichskette
{ [\epsilon]
}
{ = }{ (P(0), u(0), P'(0), u'(0))
}
{ \in }{ TTY
}
{    }{
}
{   }{
}
}
{}{}{.}
}{Unter der
\definitionsverweis {Tangentialabbildung}{}{}
\maabbdisp {T_{(P,u)} (\pi)} {T_{ (P,u)} TY} { T_PY
} {}
zu $\pi$ wird die Klasse
\mathl{[\epsilon]}{} aus (1) auf

\mavergleichskettedisp
{\vergleichskette
{ [ \pi \circ \epsilon ]
}
{ =} { [P'(0)]
}
{ } {
}
{ } {
}
{ } {
}
}
{}{}{}
abgebildet.
}{Unter der
\definitionsverweis {Tangentialabbildung}{}{}
\maabb {T(\pi)} {TTY} {TY
} {}
zu $\pi$ wird
\mathl{(P,u,v,w)}{} auf
\mathl{(P,v)}{} abgebildet.
}

}
{} {}

\inputaufgabe
{}
{

Bestimme für den Einheitskreis

\mavergleichskettedisp
{\vergleichskette
{ Y
}
{ =} { \left\{ (x_1,x_2)  \mid   x_1^2 + x_2^2 = 1         \right\}
}
{ \subset} { \R^2
}
{ } {
}
{ } {
}
}
{}{}{}
eine implizite Beschreibung des
\definitionsverweis {Tangentialbündels}{}{}
$TY$ und des zweiten Tangentialbündels $TTY$ im Sinne von
[[Differenzierbare Hyperfläche/Zweites Tangentialbündel/Induzierte Bündelhomomorphismen/Bemerkung|Bemerkung 24.5]].

}
{} {}

\inputaufgabe
{}
{

Es sei

\mavergleichskette
{\vergleichskette
{ W
}
{ \subseteq }{ \R^n
}
{  }{
}
{    }{
}
{   }{
}
}
{}{}{}
\definitionsverweis {offen}{}{,}
\maabb {\varphi} { W } { \R^k
} {}
eine
$C^2$-\definitionsverweis {stetig differenzierbare}{}{}
Abbildung,

\mavergleichskette
{\vergleichskette
{ c
}
{ \in }{ \R^k
}
{  }{
}
{    }{
}
{   }{
}
}
{}{}{}
und

\mavergleichskette
{\vergleichskette
{ Y
}
{ = }{ \varphi^{-1} (c)
}
{  }{
}
{    }{
}
{   }{
}
}
{}{}{}
die
\definitionsverweis {Faser}{}{}
über $c$, die in jedem Punkt
\definitionsverweis {regulär}{}{}
sei. Man entwickle eine implizite Beschreibung
\zusatzklammer {also mit Gleichungen} {} {}
des
\definitionsverweis {Tangentialbündels}{}{}
$TY$ und des zweiten Tangentialbündels $TTY$ ähnlich zu
[[Differenzierbare Hyperfläche/Zweites Tangentialbündel/Induzierte Bündelhomomorphismen/Bemerkung|Bemerkung 24.5]].

}
{} {}

\inputaufgabe
{}
{

Es sei $X$ eine
$C^2$-\definitionsverweis {differenzierbare Mannigfaltigkeit}{}{.}
Es seien

\mavergleichskette
{\vergleichskette
{ I,J
}
{ \subseteq }{ \R
}
{  }{
}
{    }{
}
{   }{
}
}
{}{}{}
\definitionsverweis {offene Intervalle}{}{}
mit

\mavergleichskette
{\vergleichskette
{ 0
}
{ \in }{ I,J
}
{  }{
}
{    }{
}
{   }{
}
}
{}{}{}
und es sei
\maabbeledisp {\varphi} {I \times J} {X
} {(s,t)} { \varphi(s,t)
} {,}
eine zweifach
\definitionsverweis {stetig differenzierbare}{}{}
Abbildung mit

\mavergleichskette
{\vergleichskette
{ \varphi(0,0)
}
{ = }{ P
}
{ \in }{ X
}
{    }{
}
{   }{
}
}
{}{}{.}
\aufzaehlungvier{Zeige, dass für jedes

\mavergleichskette
{\vergleichskette
{ s
}
{ \in }{ I
}
{  }{
}
{    }{
}
{   }{
}
}
{}{}{}
die Abbildung

\mathl{\varphi \vert_{ \{s\} \times J }}{} eine differenzierbare Kurve in $X$ definiert.
}{Zeige, dass durch (1) eine differenzierbare Kurve
\maabbeledisp {} { I } { TX
} { s } { [ \varphi \vert_{ \{s\} \times J } ]
} {,}
definiert wird, wobei
\mathl{[ \varphi \vert_{ \{s\} \times J } ]}{} den Tangentialvektor dieses Weges im Punkt $(s,0)$ bezeichnet.
}{Zeige, dass durch (2) ein Element in
\mathl{TTY}{} bestimmt ist.
}{Zeige, dass man jedes Element aus
\mathl{TTY}{} durch eine Abbildung $\varphi$ wie oben realisieren kann.
}

}
{} {}

\inputaufgabe
{}
{

Es sei $X$ eine
$C^2$-\definitionsverweis {differenzierbare Mannigfaltigkeit}{}{.}
Wir betrachten zweifach
\definitionsverweis {stetig differenzierbare}{}{}
Abbildungen
\maabbeledisp {\varphi} {I \times J} {X
} {(s,t)} { \varphi(s,t)
} {,}
als Elemente in
\mathl{TTX}{} im Sinne von
[[Differenzierbare_Mannigfaltigkeit/Quader/Abbildung/Zweites_Tangentialbündel/Aufgabe|Aufgabe 24.5]].
Wie verhalten sich diese Realisierungen unter der kurzen exakten Sequenz

\mathdisp {0 \, \stackrel{  }{\longrightarrow} \, p^*TX  \, \stackrel{  }{ \longrightarrow}   \, TTX  \, \stackrel{  }{\longrightarrow} \, p^*TX \,\stackrel{  } {\longrightarrow}  \, 0} { . }

}
{} {}

\inputaufgabe
{}
{

Es sei
\maabb {} { E } { X
} {}
ein
\definitionsverweis {differenzierbares Vektorbündel}{}{}
über einer
\definitionsverweis {differenzierbaren Mannigfaltigkeit}{}{}
$X$. Bestimme die
\definitionsverweis {Ränge}{}{}
für die Vektorbündel über $E$  in der kurzen exakten Sequenz

\mathdisp {0 \, \stackrel{  }{\longrightarrow} \,  V   \, \stackrel{  }{ \longrightarrow}   \,  TE  \, \stackrel{  }{\longrightarrow} \, p^* TX \,\stackrel{  } {\longrightarrow}  \, 0} { . }

}
{} {}

\inputaufgabe
{}
{

Bestimme die
\definitionsverweis {Zusammenhänge}{}{}
auf dem Nullbündel auf einer
\definitionsverweis {differenzierbaren Mannigfaltigkeit}{}{}
$X$.

}
{} {}

\inputaufgabe
{}
{

Bestimme die
\definitionsverweis {Zusammenhänge}{}{}
auf einem
\definitionsverweis {Vektorbündel}{}{}
auf einer einpunktigen
\definitionsverweis {Mannigfaltigkeit}{}{}
$X$.

}
{} {}

\inputaufgabe
{}
{

Es sei $X$ eine
\definitionsverweis {differenzierbare Mannigfaltigkeit}{}{}
und $E$ ein
\definitionsverweis {differenzierbares Vektorbündel}{}{}
auf $X$, das mit einem
\definitionsverweis {Zusammenhang}{}{}
versehen sei. Zeige, dass man das Konzept vertikale Ableitung auch längs einer differenzierbaren Abbildung
\maabbdisp {\psi} { Z } { X
} {}
einführen kann, vergleiche hierzu auch
[[Hyperfläche/Tangentiales Vektorfeld/Längs einer Kurve/Kovariante Ableitung/Definition|Definition 6.5]].

}
{} {}

  \zwischenueberschrift{Aufgaben zum Abgeben}

\inputaufgabe
{3}
{

Es sei $X$ ein
\definitionsverweis {topologischer Raum}{}{}
und sei

\mathdisp {0 \, \stackrel{  }{\longrightarrow} \, U  \, \stackrel{  }{ \longrightarrow}   \, V  \, \stackrel{  }{\longrightarrow} \, W \,\stackrel{  } {\longrightarrow}  \, 0} {  }

eine
\definitionsverweis {kurze exakte Sequenz}{}{}
von
\definitionsverweis {Vektorbündeln}{}{}
über $X$. Es sei
\maabb {\varphi} {Y} {X
} {}
eine
\definitionsverweis {stetige Abbildung}{}{}
von einem weiteren topologischen Raum $Y$. Zeige, dass der Rückzug zu einer kurzen exakten Sequenz

\mathdisp {0 \, \stackrel{  }{\longrightarrow} \, \varphi^*U  \, \stackrel{  }{ \longrightarrow}   \, \varphi^*V  \, \stackrel{  }{\longrightarrow} \, \varphi^*W \,\stackrel{  } {\longrightarrow}  \, 0} {  }

von Vektorbündeln auf $Y$ führt.

}
{} {}

\inputaufgabe
{3 (1+2)}
{

Es sei

\mavergleichskette
{\vergleichskette
{ I
}
{ \subseteq }{ \R
}
{  }{
}
{    }{
}
{   }{
}
}
{}{}{}
ein
\definitionsverweis {offenes Intervall}{}{,}
wir betrachten die
\definitionsverweis {kurze exakte Sequenz}{}{}
\zusatzklammer {von Vektorbündeln über $I \times \R$} {} {}

\mathdisp {0 \, \stackrel{  }{\longrightarrow} \, I \times \R \times \R   \, \stackrel{  }{ \longrightarrow}   \,  I \times \R \times \R \times \R   \, \stackrel{  }{\longrightarrow} \,  I \times \R \times \R  \,\stackrel{  } {\longrightarrow}  \, 0} { , }

wobei die erste Abbildung durch
\mathl{(P,u,w) \mapsto (P,u,0,w)}{} und die zweite Abbildung durch
\mathl{(P,u,v,w) \mapsto (P,u,v)}{} gegeben sei. Es seien
\maabbdisp {g,h} {I \times \R } {\R
} {}
stetig differenzierbare Funktionen.
\aufzaehlungdrei{Zeige, dass durch
\maabbeledisp {} {I \times \R \times \R \times \R  } { I \times \R \times \R
} {(P,u,v,w)} { (P,u,   v g(P,u) +w h(P,u) )
} {,}
ein Vektorbündelhomomorphismus über $I \times \R$ gegeben ist.
}{Es sei

\mavergleichskette
{\vergleichskette
{ h
}
{ = }{ 1
}
{  }{
}
{    }{
}
{   }{
}
}
{}{}{}
konstant. Zeige, dass die Abbildung in (1) eine Spaltung des Bündels in der Mitte definiert.
}{
}

}
{} {}

\inputaufgabe
{4}
{

Bestimme für die ebene Kurve

\mavergleichskettedisp
{\vergleichskette
{ Y
}
{ =} { \left\{ (x_1,x_2)  \mid   x_1^3 -5x_1x_2 + x_2^4 = 1         \right\}
}
{ \subset} { \R^2
}
{ } {
}
{ } {
}
}
{}{}{}
eine implizite Beschreibung des
\definitionsverweis {Tangentialbündels}{}{}
$TY$ und des zweiten Tangentialbündels $TTY$ im Sinne von
[[Differenzierbare Hyperfläche/Zweites Tangentialbündel/Induzierte Bündelhomomorphismen/Bemerkung|Bemerkung 24.5]].

}
{} {}

\inputaufgabe
{5 (3+2)}
{

Wir betrachten die differenzierbare Fläche

\mavergleichskettedisp
{\vergleichskette
{ Y
}
{ =} { \left\{ (x_1,x_2,x_3)  \mid   x_1^2+x_2^2 -x_3^2 = 0  , \, (x,y,z) \neq 0       \right\}
}
{ \subseteq} { W
}
{ =} { \R^3 \setminus \{(0,0,0) \}
}
{ \subset} { \R^3
}
}
{}{}{.}
\aufzaehlungzwei {Bestimme eine implizite Beschreibung des
\definitionsverweis {Tangentialbündels}{}{}
$TY$ und des zweiten Tangentialbündels $TTY$ im Sinne von
[[Differenzierbare Hyperfläche/Zweites Tangentialbündel/Induzierte Bündelhomomorphismen/Bemerkung|Bemerkung 24.5]].
} {Berechne für die differenzierbare Kurve
\maabbeledisp {} { I } { Y
} { t } { \left(  \cos  t   , \,   \sin  t  , \,   1                 \right)
} {}
die zughörige Kurve in das Tangentialbündel und in das zweite Tangentialbündel.
}

}
{} {}

\inputaufgabe
{3}
{

Es sei

\mavergleichskette
{\vergleichskette
{ U
}
{ \subseteq }{ \R^n
}
{  }{
}
{    }{
}
{   }{
}
}
{}{}{}
\definitionsverweis {offen}{}{}
und sei

\mavergleichskette
{\vergleichskette
{ E
}
{ = }{ U \times \R
}
{  }{
}
{    }{
}
{   }{
}
}
{}{}{,}
wobei wir Funktionen
\maabb {f} { U } {\R
} {}
als
\definitionsverweis {Schnitte}{}{}
in $E$ auffassen. Es sei $E$ mit dem
\definitionsverweis {trivialen Zusammenhang}{}{}
versehen. Zeige

\mavergleichskette
{\vergleichskette
{ \nabla (f)
}
{ = }{ df
}
{  }{
}
{    }{
}
{   }{
}
}
{}{}{}
für stetig differenzierbare Funktionen.

}
{} {}

 [[Kategorie:Latexseite]]
